\section{Score Pretraining and Joint Receiver Training}
\label{sec:learning}

Training the deep-unfolded MP receiver in Section~\ref{sec:algorithm}
starts with learning a score-based channel prior. Score pretraining
uses Gaussian-perturbed samples of multiantenna AFDM effective
channels generated with continuous path parameters. The pretrained
network initializes the channel denoisers. We then jointly optimize
the denoisers, activity classifiers, and message parameters using
received pilot observations, allowing channel and activity errors
to guide the updates across receiver layers.

\subsection{First-Order Score Matching}

Let $h_R\in\mathbb R^{2D}$ denote a normalized channel sample.
We perturb it as $u=h_R+\sqrt{v/2}\,\varepsilon$, where
$\varepsilon\sim\mathcal N(0,\I_{2D})$. Denote the perturbed
density by $q_{1,R}(u;v)$ and its score by
$s^\star(u,v)=\nabla_u\log q_{1,R}(u;v)$.
Tweedie's formula relates this score to the conditional channel mean
\cite{efron2011,cai2026}:
\begin{equation}
 \E[h_R\mid u]=u+\frac v2\,s^\star(u,v).
 \label{eq:twc-tweedie}
\end{equation}
Equation~\eqref{eq:twc-tweedie} expresses the conditional mean in
\eqref{eq:twc-local-channel-target} through the score of the perturbed
channel distribution. Replacing this score by a network gives the initial denoising
correction in \eqref{eq:twc-conditional-mean} with $a_\ell=1$.
Joint receiver training subsequently optimizes the network and this
coefficient for channel estimation through the unfolded layers.

To learn the score over the input noise levels, consider the
weighted score matching objective \cite{cai2026}
\begin{equation}
 \mathcal L_{\rm SM}(\theta)
 =\E_{v,u}\!\left[
 \frac{v}{4D}\left\|s_\theta(u,v)
       -\nabla_u\log q_{1,R}(u;v)\right\|^2\right],
 \label{eq:twc-score-matching}
\end{equation}
where $u$ follows $q_{1,R}(\cdot;v)$ at each noise level.
The marginal score in this objective is unknown. Denoising score
matching replaces it by the conditional perturbation score and has
the same minimizers \cite{vincent2011,song2019}.
The Gaussian perturbation has conditional score
$\nabla_u\log q(u\mid h_R,v)=-\varepsilon/\sqrt{v/2}$.
Substitution gives the training loss
\begin{equation}
 \mathcal L_{\rm DSM}(\theta)=\E_{h_R,v,\varepsilon}
 \left[\frac{\|\sqrt{v/2}\,s_\theta(h_R+\sqrt{v/2}\varepsilon,v)
                      +\varepsilon\|^2}{2D}\right].
 \label{eq:twc-dsm}
\end{equation}
Training samples are drawn from the channel distribution before
activity masking and perturbed over a range of noise levels.
The resulting network initializes $s_{\theta_\ell}$ in
\eqref{eq:twc-conditional-mean}. During joint receiver training, this
score-based form is retained while its network weights and denoising
coefficient are optimized for the inputs produced by the unfolded
layers.

\subsection{Network Inputs}

For each terminal, the real and imaginary antenna maps are arranged
as $2M$ channels over the $N\times N$ DAFT coordinates. The score
network is conditioned on $v_\ell$ and returns a tensor of the same
shape, which is flattened to the vector used in
\eqref{eq:twc-conditional-mean}. This representation retains the full
multiantenna channel within each network evaluation.

The activity classifier concatenates the denoiser input and its
channel estimate into $4M$ input channels and scales them by
$(1+v_\ell)^{-1/2}$. Three convolutional layers followed by spatial
average pooling extract a feature vector. Two fully connected layers
combine this vector with the five scalar inputs
\begin{equation}
 \begin{aligned}
 \Big(&\log v_\ell,\
 \log\!\left(1+\frac{\|u_{k\ell}\|^2}{Dv_\ell}\right),\
 \log\!\left(1+\frac{\|\mu_{k\ell}\|^2}{Dv_\ell}\right),\\
 &\log\!\left(1+\frac{\|u_{k\ell}-\mu_{k\ell}\|^2}{Dv_\ell}\right),\
 \logit\rho\Big)
 \end{aligned}
 \label{eq:twc-classifier-features}
\end{equation}
to produce the log-odds correction in \eqref{eq:twc-activity}.
The energy terms describe the denoiser input, channel estimate, and denoising
residual relative to the input noise level. Each network acts on
one terminal at a time, with parameters shared across terminals
within a layer. Thus its input dimensions depend on $M$ and $N$,
while $K$ determines the number of evaluations and $Q$ enters through
the linear estimator.

\subsection{Joint Channel and Activity Training}

To adapt the denoisers to the inputs produced by the receiver, we
pass each training observation through the updates in
Section~\ref{sec:algorithm}. In this training stage, the sensing matrix
and noise variance are supplied from the physical model, as specified
in Section~\ref{sec:system}; phase and noise estimation are performed
separately in the complete receiver.
The denoiser inputs and feedback messages depend on the current layer
parameters. We supervise the two factors in
\eqref{eq:twc-mixture-mean} using the squared channel error on active
users and binary cross-entropy on all users. These losses provide
distinct supervision for the conditional channel estimate and activity
probability within a common objective. Their
product determines the feedback to later layers, so both estimators
participate in joint optimization.

Suppose the receiver training set contains $B$ scenes, with $K_i$
candidate users in scene $i$. Let $h_{ik,R}$ and $\alpha_{ik}$
denote the channel and activity label of user $k$ in scene $i$.
Training and validation use disjoint physical scenes.
The channel and activity losses at layer $\ell$ are respectively
given by
\begin{align}
 R_H^{(\ell)}&=\frac{1}{E_{\rm tr}}\sum_{i,k}
       \alpha_{ik}\|\mu_{ik\ell}-h_{ik,R}\|^2,
 \label{eq:twc-channel-loss}\\
 R_a^{(\ell)}&=-\frac{1}{T_{\rm tr}}\sum_{i,k}
       \left[\alpha_{ik}\log p_{ik\ell}
       +(1-\alpha_{ik})\log(1-p_{ik\ell})\right].
 \label{eq:twc-activity-loss}
\end{align}
Here $E_{\rm tr}=\sum_{i,k}\alpha_{ik}\|h_{ik,R}\|^2$ and
$T_{\rm tr}=\sum_i K_i$ are the fixed active-channel energy and user
count of the training set. The outputs $\mu_{ik\ell}$ and $p_{ik\ell}$
are computed from the received observations; the activity labels
select the channel terms in \eqref{eq:twc-channel-loss}.

For a receiver with $m$ layers, the joint training objective is
\begin{equation}
 J_m=\frac{\sum_{\ell=1}^{m}\omega_\ell
                  \left(R_H^{(\ell)}+R_a^{(\ell)}\right)}
                 {\sum_{\ell=1}^{m}\omega_\ell},
 \label{eq:twc-layer-objective}
\end{equation}
where $\omega_\ell$ is the weight assigned to layer $\ell$.
The last layer provides the final estimates, while earlier-layer
outputs provide intermediate supervision.
Gradients of \eqref{eq:twc-layer-objective} are propagated through
the denoisers, classifiers, and inter-layer messages. Thus activity
supervision also adjusts channel denoising, and channel losses at later
layers refine the activity estimates used in earlier feedback.

\subsection{Progressive Training}

We increase receiver depth progressively. Each new layer has independent
parameters, with its score network initialized from the same
pretrained network.
Its classifier output $g_{\phi_\ell}$ and the parameters $c_{\alpha,\ell}$,
$c_{\beta,\ell}$, and $c_{\gamma,\ell}$ are initialized to zero,
so the initial activity probability equals $\rho$.
Earlier layers retain their previously selected parameters.

At each depth, we first train the new layer and adjust the variance
scaling factor used in the preceding layer's outgoing message. We then
fine-tune all layers jointly using a reduced learning rate. Both stages minimize
\eqref{eq:twc-layer-objective}. The checkpoint with the smallest
final-layer channel and activity loss on the validation set initializes
the next depth. Validation uses the training normalization constants.
After training to depth $L$, the selected layer parameters are fixed
for the receiver summarized in Section~\ref{sec:overall}.
